\documentclass[11pt,a4paper]{article}

% ── Codificación y fuentes ────────────────────────────────────────────────────
\usepackage[utf8]{inputenc}
\usepackage[T1]{fontenc}
\usepackage{lmodern}
\usepackage[spanish]{babel}

% ── Geometría y espaciado ─────────────────────────────────────────────────────
\usepackage[top=2.5cm, bottom=2.5cm, left=2.8cm, right=2.8cm]{geometry}
\usepackage{setspace}
\setstretch{1.15}
\usepackage{parskip}

% ── Tablas ────────────────────────────────────────────────────────────────────
\usepackage{booktabs}
\usepackage{tabularx}
\usepackage{array}
\usepackage{longtable}
\usepackage{enumitem}

% ── Colores y cajas ───────────────────────────────────────────────────────────
\usepackage[dvipsnames]{xcolor}
\usepackage{tcolorbox}
\tcbuselibrary{skins, breakable}

% ── Cabeceras y pies de página ────────────────────────────────────────────────
\usepackage{fancyhdr}
\usepackage{lastpage}
\pagestyle{fancy}
\fancyhf{}
\fancyhead[L]{\small\textbf{Informe de Evaluación} --- Física para Bioanalistas}
\fancyhead[R]{\small Deliangelis Inojosa}
\fancyfoot[C]{\small Página \thepage\ de \pageref{LastPage}}
\renewcommand{\headrulewidth}{0.4pt}

% ── Hiperenlaces ──────────────────────────────────────────────────────────────
\usepackage[colorlinks=true, linkcolor=NavyBlue, urlcolor=NavyBlue]{hyperref}

% ── Paleta de colores personalizados ─────────────────────────────────────────
\definecolor{desempeno_excelente}{HTML}{1A6B2F}
\definecolor{desempeno_bueno}{HTML}{2E5A9C}
\definecolor{desempeno_suficiente}{HTML}{8B6914}
\definecolor{desempeno_deficiente}{HTML}{C0392B}
\definecolor{desempeno_insuficiente}{HTML}{7B241C}
\definecolor{grisFondo}{HTML}{F5F5F5}
\definecolor{azulTitulo}{HTML}{1B2A6B}
\definecolor{verdeFortaleza}{HTML}{1A6B2F}
\definecolor{naranjaMejora}{HTML}{A04000}

% ── Estilos de cajas ─────────────────────────────────────────────────────────
\tcbset{
  cajaTitulo/.style={
    enhanced, breakable,
    colback=azulTitulo!8, colframe=azulTitulo,
    fonttitle=\bfseries\large, coltitle=white,
    attach boxed title to top left={yshift=-2mm, xshift=4mm},
    boxed title style={colback=azulTitulo},
    top=6pt, bottom=6pt, left=8pt, right=8pt,
  },
  cajaFortaleza/.style={
    enhanced, breakable,
    colback=verdeFortaleza!6, colframe=verdeFortaleza!60,
    fonttitle=\bfseries, coltitle=verdeFortaleza,
    top=4pt, bottom=4pt, left=8pt, right=8pt,
  },
  cajaMejora/.style={
    enhanced, breakable,
    colback=naranjaMejora!6, colframe=naranjaMejora!60,
    fonttitle=\bfseries, coltitle=naranjaMejora,
    top=4pt, bottom=4pt, left=8pt, right=8pt,
  },
  cajaRecomendacion/.style={
    enhanced, breakable,
    colback=grisFondo, colframe=gray!50,
    fonttitle=\bfseries, coltitle=black,
    top=4pt, bottom=4pt, left=8pt, right=8pt,
  },
}

% ─────────────────────────────────────────────────────────────────────────────
\begin{document}

% ══ PORTADA ══════════════════════════════════════════════════════════════════
\begin{center}
  {\Large\bfseries\color{azulTitulo} Universidad de Oriente/Núcleo Sucre --- Física para Bioanalistas}\\[4pt]
  {\large Informe de Evaluación Preliminar}\\[12pt]
  \rule{\linewidth}{1.2pt}\\[6pt]
  {\LARGE\bfseries Deliangelis Inojosa}\\[4pt]
  \rule{\linewidth}{0.6pt}
\end{center}

% ══ FICHA TÉCNICA ════════════════════════════════════════════════════════════
\vspace{4pt}
\begin{tcolorbox}[cajaTitulo, title=Datos del informe]
\begin{tabular}{@{}llll@{}}
  \textbf{Estudiante:}  & Deliangelis Inojosa  &
  \textbf{Fecha:}       & 2026-08-08 \\[4pt]
  \textbf{Archivo PDF:} & \texttt{Deliangelis\_Inojosa.pdf}  &
  \textbf{Páginas:}     & 5 \\
\end{tabular}
\end{tcolorbox}

% ══ RESULTADO GLOBAL ═════════════════════════════════════════════════════════
\vspace{6pt}
\begin{center}
  \begin{tcolorbox}[
    enhanced,
    colback=white, colframe=desempeno_bueno,
    width=0.62\linewidth,
    boxrule=2pt,
    halign=center, valign=center,
    top=8pt, bottom=8pt,
  ]
    \centering
    {\large\bfseries Puntaje sugerido}\\[6pt]
    {\Huge\bfseries\color{desempeno_bueno} 7.5/10}\\[6pt]
    {\large Nivel de desempeño: \textbf{\color{desempeno_bueno} Bueno}}
  \end{tcolorbox}
\end{center}

% ══ RESUMEN GENERAL ══════════════════════════════════════════════════════════
\section*{Resumen general}
El estudiante demuestra una buena comprensión de los principios físicos fundamentales relacionados con los líquidos y fenómenos de transporte, especialmente en mecánica respiratoria, termorregulación y capilaridad. Sin embargo, se observan deficiencias recurrentes en la ejecución de cálculos numéricos, donde a menudo se omiten pasos finales o se utilizan valores incorrectos para constantes. La argumentación cualitativa es generalmente correcta, pero en ocasiones podría ser más profunda y explícita.

% ══ FORTALEZAS Y ÁREAS DE MEJORA ════════════════════════════════════════════
\vspace{4pt}
\begin{minipage}[t]{0.48\linewidth}
  \begin{tcolorbox}[cajaFortaleza, title={Fortalezas}]
\begin{itemize}[leftmargin=1.5em, itemsep=2pt]
  \item Buena comprensión conceptual de la Ley de Laplace y el rol del surfactante en la mecánica respiratoria.
  \item Correcta interpretación de los efectos osmóticos (soluciones hipotónicas e hipertónicas) sobre los eritrocitos.
  \item Correcta aplicación cualitativa de la Ley de Jurin para explicar el ascenso capilar y el efecto de la hidrofobicidad.
  \item Identificación correcta de los riesgos clínicos asociados a los fenómenos físicos (golpe de calor, ruptura de membrana, hemólisis, crenación).
  \item Manejo adecuado de unidades en los cálculos completados.
\end{itemize}
  \end{tcolorbox}
\end{minipage}
\hfill
\begin{minipage}[t]{0.48\linewidth}
  \begin{tcolorbox}[cajaMejora, title={Areas de mejora}]
\begin{itemize}[leftmargin=1.5em, itemsep=2pt]
  \item Completar todos los cálculos numéricos hasta obtener un resultado final con sus unidades correctas.
  \item Verificar la correcta sustitución de todas las constantes y valores en las fórmulas, prestando atención a la notación científica.
  \item Profundizar en las explicaciones cualitativas, justificando los principios físicos subyacentes de manera más explícita y detallada.
  \item Mejorar la claridad en la escritura de valores numéricos en las sustituciones.
\end{itemize}
  \end{tcolorbox}
\end{minipage}

% ══ OBSERVACIONES POR PREGUNTA ═══════════════════════════════════════════════
\section*{Observaciones por pregunta}

\begin{longtable}{>{\bfseries}p{2.8cm} p{1.6cm} p{7.0cm} p{2.8cm}}
  \toprule
  \textbf{Pregunta} & \textbf{Puntaje} & \textbf{Evaluación} & \textbf{Errores detectados} \\
  \midrule
  \endhead
  \bottomrule
  \endfoot
    \textbf{Pregunta 1: Termorregulación y Eficiencia de la Sudoración} & 1.5/2.0 & La respuesta a la parte a) es correcta en su conclusión, pero carece de la explicación fundamental sobre el calor latente de vaporización. La parte b) es concisa y correcta. La parte c) presenta un cálculo numérico correcto, incluyendo la conversión de unidades. & \textit{Falta de explicación detallada sobre el principio físico (calor latente de vaporización) en la parte a).} \\
    \midrule
    \textbf{Pregunta 2: Mecánica Respiratoria y Ley de Laplace} & 2.0/2.0 & Todas las partes de esta pregunta son respondidas de manera correcta y clara. La explicación de la Ley de Laplace y el rol del surfactante es precisa, y el cálculo numérico es exacto con unidades correctas. & \textit{---} \\
    \midrule
    \textbf{Pregunta 3: Optimización en Hemodiálisis y Primera Ley de Fick} & 1.5/2.0 & La parte a) demuestra una correcta aplicación de la Ley de Fick y el cálculo del aumento de eficiencia. La parte b) identifica correctamente el riesgo principal. Sin embargo, en la parte c), aunque la fórmula es correcta, la sustitución del gradiente de concentración es incompleta/ambigua y el cálculo final no se presenta. & \textit{Cálculo incompleto en la parte c).; Sustitución ambigua o incorrecta del gradiente de concentración en la parte c).} \\
    \midrule
    \textbf{Pregunta 4: Errores de Infusión Intravenosa y Ósmosis} & 1.0/2.0 & Las partes a) y b) son conceptualmente correctas y explican adecuadamente los fenómenos de hemólisis y crenación. En la parte c), la fórmula de presión osmótica es correcta, pero los valores sustituidos para la constante de los gases (R) y la temperatura (T) son incorrectos o están incompletos, y el cálculo final no se realiza. & \textit{Sustitución incorrecta o incompleta de las constantes R y T en la fórmula de presión osmótica en la parte c).; Cálculo incompleto en la parte c).} \\
    \midrule
    \textbf{Pregunta 5: Capilaridad y Toma de Muestras Sanguíneas} & 1.5/2.0 & Las partes a) y b) son conceptualmente correctas, explicando el efecto de la hidrofobicidad y la relación inversa con el radio. En la parte c), la fórmula de la Ley de Jurin es correcta, pero la sustitución del radio 'r' es ambigua ('2.0 X 1' en lugar de '2.0x10\textasciicircum{}-4') y el cálculo final no se completa. & \textit{Sustitución ambigua del radio 'r' en la parte c).; Cálculo incompleto en la parte c).} \\
    \midrule
\end{longtable}

% ══ ERRORES DETECTADOS ═══════════════════════════════════════════════════════
\section*{Análisis de errores}

\subsection*{Errores conceptuales}
\begin{itemize}[leftmargin=1.5em, itemsep=2pt]
  \item En la Pregunta 1a, la explicación sobre la ineficiencia del sudor que gotea podría haber sido más explícita sobre la falta de aprovechamiento del calor latente de vaporización, lo cual es un punto clave en la termorregulación por sudoración.
\end{itemize}

\subsection*{Errores procedimentales}
\begin{itemize}[leftmargin=1.5em, itemsep=2pt]
  \item Cálculo incompleto en la Pregunta 3c, con una sustitución poco clara del gradiente de concentración.
  \item Sustitución incorrecta o incompleta de las constantes R y T en la fórmula de presión osmótica en la Pregunta 4c, resultando en un cálculo incompleto.
  \item Sustitución ambigua del valor del radio 'r' en la Pregunta 5c, y cálculo incompleto.
\end{itemize}

% ══ RECOMENDACIONES AL ESTUDIANTE ════════════════════════════════════════════
\section*{Retroalimentación para el estudiante}
\begin{tcolorbox}[cajaRecomendacion, title={Dirigido a: Deliangelis Inojosa}]
Estimado estudiante, has demostrado una sólida base conceptual en varios temas de la unidad, especialmente en la mecánica respiratoria, los efectos osmóticos y la capilaridad. Tus explicaciones cualitativas son, en su mayoría, correctas y pertinentes. Sin embargo, es crucial que prestes más atención a la sección de cálculos. Asegúrate de sustituir correctamente todos los valores, incluyendo las constantes, y de completar cada cálculo hasta obtener un resultado numérico final con sus unidades apropiadas. Practicar más la resolución completa de problemas numéricos te ayudará a consolidar tu comprensión y a evitar errores por omisión o sustitución incorrecta. Además, intenta expandir un poco más tus justificaciones cualitativas, explicando el 'porqué' de los fenómenos con mayor detalle físico.
\end{tcolorbox}

% ══ NOTA PARA EL PROFESOR ════════════════════════════════════════════════════
\section*{Nota para el profesor}
\begin{tcolorbox}[
  enhanced, breakable,
  colback=yellow!8, colframe=yellow!60!black,
  fonttitle=\bfseries, coltitle=black,
  title={Observaciones para revisión manual},
  top=4pt, bottom=4pt, left=8pt, right=8pt,
]
El estudiante presenta un buen nivel de comprensión conceptual general. Los errores principales se concentran en la fase de cálculo numérico: sustituciones incompletas o incorrectas de constantes (ej. R y T en ósmosis), notación ambigua para valores (ej. radio en capilaridad), y la omisión de la finalización de los cálculos. Las explicaciones cualitativas son correctas pero a veces concisas, pudiendo beneficiarse de una mayor profundidad en la justificación de los principios físicos. El puntaje sugerido refleja esta mezcla de buena comprensión conceptual y deficiencias en la ejecución numérica.
\end{tcolorbox}

% ══ PIE DE INFORME ════════════════════════════════════════════════════════════
\vspace{12pt}
\begin{center}
  \footnotesize\color{gray}
  Informe generado automáticamente por el sistema de evaluación con IA (gemini-2.5-flash) y soporte LaTex. \\
  Este documento es una evaluación \textbf{preliminar} para ser revisado por el profesor antes de ser entregado al 
  estudiante. \\
  Generado el 2026-08-08 \textbullet\ BIOANALISIS Sistema Académico de Evaluación.
  \end{center}

\end{document}
